% 英文摘要頁
\begin{EnAbstract}
    \begin{EnAbstractItems}
        % 關鍵詞，請自己填，多個關鍵字以逗號 "," 隔開
        \noindent \text Keyword: Large Language Models, Multi-Agent Systems, RESTful API, Automated Software Testing, Reinforcement Learning, Model Context Protocol, Agent-to-Agent Protocol, AI Agents, Workflow Orchestration, Software Quality Assurance

    \end{EnAbstractItems}

    \begin{EnAbstractDescription}
        \text As software development shifts toward agile methodologies and microservice architectures, RESTful APIs have become the core interface for system integration and service exchange. However, existing automated testing approaches continue to struggle with multi-step operation dependencies, semantic constraints, and dynamic real-world execution environments, resulting in insufficient semantic understanding, difficulties in cross-operation dependency reasoning, and a lack of recovery capability in testing workflows. These limitations lead to a high proportion of invalid requests and constrain overall test completion rates.
        
        \text To address these challenges, this study proposes a multi-agent RESTful API automated testing framework centered on a Post-Execution Workflow Recovery Layer. The framework employs Large Language Models (LLMs) as the semantic reasoning core, integrating AutoGen for multi-agent collaboration, LangGraph for workflow control, Model Context Protocol (MCP), and Agent-to-Agent Protocol (A2A) to establish a complete testing pipeline spanning OpenAPI specification parsing, test case generation, test code materialization, API execution, workflow decision-making, and result analysis.

        \text The primary focus of this study is not merely generating test requests, but addressing the workflow decision problem that arises after API execution. To this end, post-execution state transitions are formalized as a Markov Decision Process (MDP), enabling the system to select subsequent workflow actions based on signals including HTTP status codes, error types, dependency states, authentication states, and resource availability. Candidate actions include retry, token refresh, resource resolution, test data regeneration, and result reporting, with the goal of improving workflow continuity and recovery capability. Within this design, the workflow recovery layer provides executable recovery mechanisms for scenarios such as 404 errors, resource unavailability, and dependency failures, while the workflow decision policy is responsible for selecting the next action among candidate recovery operations.

        \text For empirical evaluation, this study uses 100 manual-doc-grounded tasks from the TMDB dataset as the primary benchmark, comparing test performance across versions v1 to v5. Results show pass rates of 0.25, 0.72, 0.69, 0.69, and 0.79 for v1 through v5 respectively. Notably, v5 maintains a Strict Pass Rate of 0.69, identical to v3 and v4, while its Final Pass Rate improves to 0.79 with 10 additional Workflow Recovered Pass cases, indicating that the primary benefit derives from post-execution workflow recovery rather than improvements in upstream test generation. Further analysis reveals that current recovery gains are concentrated in 404 resource dependency failure scenarios.

        \text Regarding offline reinforcement learning, this study constructs an offline dataset from workflow transition events and trains an Offline-Q Policy using Tabular Fitted Q Iteration. The primary purpose of this component is to explore the feasibility of formalizing the workflow decision problem as a learning-based workflow decision strategy. Results demonstrate that the learned policy successfully captures partial workflow decision patterns and provides incremental benefits in specific recovery scenarios. However, due to constraints including a single API domain, Live API environment variability, and limited offline dataset scale, the overall performance is insufficient to claim comprehensive superiority over the rule-based baseline. Overall, this study validates the feasibility of integrating large language models, multi-agent collaboration, and workflow decision strategies into RESTful API automated testing, and provides a concrete reference architecture for the design and evaluation of future intelligent API testing systems.
    \end{EnAbstractDescription}
    
\end{EnAbstract}

